\documentclass[11pt,a4paper]{article}
\usepackage[margin=1.8cm]{geometry}
\usepackage[T1]{fontenc}
\usepackage[utf8]{inputenc}
\usepackage[brazil]{babel}
\usepackage{booktabs}
\usepackage{tabularx}
\usepackage{microtype}
\usepackage{hyperref}
\hypersetup{colorlinks=true,linkcolor=black,urlcolor=blue}
\pagestyle{empty}
\setlength{\parskip}{0.35em}
\setlength{\parindent}{0pt}

\title{\vspace{-1.2cm}\textbf{Team Shannon} --- Entrega de meio de semestre\\
\large Motor de busca: Porto de Santos (Baixada Santista)}
\author{Adriane da Costa Santos · Danilo Prado de Lima Silva · Victória Cabral Quintério\\
\small Fatec Rubens Lara · PI III · Ciência de Dados\\
\url{https://github.com/danilopradooficial/projetointegrador3}}
\date{}

\begin{document}
\maketitle
\vspace{-1.2cm}

\noindent\textit{Documento de auditoria (extra). O relatório oficial da entrega é a ficha
\texttt{consolidados/00\_FICHA\_PROJETO.md}. Demonstração: \texttt{Rscript aula02.R}
e \texttt{Rscript entrega-meio-semestre-demo.R}.}

\section*{1. Por que estes dados}
Tema ligado à \textbf{Baixada Santista} via Porto de Santos. Três artigos da Wikipédia
(CC~BY-SA) cobrem faces distintas do mesmo objeto:

\begin{itemize}
  \setlength\itemsep{0.1em}
  \item \textbf{Local} --- \emph{Porto de Santos}: estuário, economia, cargas, infraestrutura.
  \item \textbf{Empresa pública} --- \emph{Autoridade Portuária de Santos} (APS / Codesp):
        quem administra, Landlord Port, lei~8.630.
  \item \textbf{Fundador} --- \emph{Francisco de Paula Ribeiro}: engenheiro e origem histórica.
\end{itemize}

Um documento de recuperação é uma \textbf{frase pontuada} ($dN.k$), não o artigo inteiro:
é a menor unidade que responde sozinha a uma consulta digitável. Os IDs são canônicos de
01a até 05b (\texttt{01a-ler-frases.R}).

\section*{2. Cara do corpus}
\begin{center}
\begin{tabular}{@{}lrrr@{}}
\toprule
Fonte & Frases & Palavras (mediana) & Papel \\
\midrule
Porto de Santos ($d1$) & 88 & 26 & lugar / economia \\
Autoridade Portuária ($d2$) & 29 & 26 & gestão pública \\
F.\ de Paula Ribeiro ($d3$) & 5 & --- & fundador \\
\midrule
\textbf{Total} & \textbf{122} & média $\approx$ 28,6 & vocab.\ stems $\approx$ 873 \\
\bottomrule
\end{tabular}
\end{center}
\noindent\textit{Nota:} o PDF da entrega cita 20--60 parágrafos; estamos acima porque a Aula~05
já fatiou o mesmo corpus em frases --- sem páginas novas.

\section*{3. Três consultas de trabalho}
\begin{enumerate}
  \setlength\itemsep{0.15em}
  \item \texttt{localizacao porto santos guarujá cubatão} --- lugar.
  \item \texttt{quem administra porto santos autoridade} --- autoridade pública.
  \item \texttt{francisco de paula ribeiro porto} --- fundador.
\end{enumerate}

\section*{4. Auditoria: mesma pergunta, vários métodos}
Tempos em ms (Python de referência; R no script de demo). Qualidade: $P@3$ / nDCG com
gabarito da Aula~05a (limiar grau $\geq 2$); ``---'' = sem relevantes nesse limiar.

\subsection*{Consulta 1 --- localização (\texttt{qrels} q15, $R{=}3$)}
\begin{center}
\small
\begin{tabular}{@{}l r l r r@{}}
\toprule
Método & ms & Top-3 & $P@3$ & nDCG \\
\midrule
AND & 1,5 & \textit{(vazio --- exige todos os stems)} & 0 & 0 \\
OR & 0,9 & conjunto ($\lvert$OR$\rvert{=}60$), sem ordem & 0,33 & 1,00 \\
TF-IDF & 2,1 & d1.1 · d1.61 · d1.66 & 0,67 & 1,00 \\
Cosseno & 234 & d1.1 · d1.61 · d1.30 & 0,67 & 0,97 \\
BM25 & 2,1 & d1.1 · d1.61 · d1.82 & 0,67 & 1,00 \\
\bottomrule
\end{tabular}
\end{center}
\noindent 1\textsuperscript{o} cosseno/BM25: \emph{``Porto de Santos é um porto estuarino, localizado
nos municípios de Santos, Guarujá e Cubatão\ldots''} --- resposta direta.

\subsection*{Consulta 2 --- quem administra (\texttt{qrels} q02, $R{=}0$ grau$\geq$2)}
\begin{center}
\small
\begin{tabular}{@{}l r l@{}}
\toprule
Método & ms & Top-3 \\
\midrule
AND & 0,9 & vazio \\
TF-IDF & 1,8 & d1.66 · d1.13 · d2.16 \\
Cosseno / BM25 & 301 / 2,3 & d1.36 · d2.9 · d2.4 \\
\bottomrule
\end{tabular}
\end{center}
\noindent Pool com poucos grau~2 nesta necessidade: métricas binárias não discriminam.
Texto do topo fala de áreas administrativas do Porto Organizado --- qualidade
parcial (gestão), não ``quem é a APS'' em uma linha.

\subsection*{Consulta 3 --- fundador (\texttt{qrels} q04, $R{=}1$)}
\begin{center}
\small
\begin{tabular}{@{}l r l r@{}}
\toprule
Método & ms & Top-3 & $P@3$ \\
\midrule
AND & 0,7 & vazio & 0 \\
TF-IDF / Cosseno / BM25 & 3--5 / 347 / 5 & \textbf{d3.1} · d1.17 · d3.4/d3.5 & 0,33 \\
\bottomrule
\end{tabular}
\end{center}
\noindent 1\textsuperscript{o}: biografia de Francisco de Paula Ribeiro ($d3.1$) --- acerto imediato
(MRR~1 nos ranqueadores).

\section*{5. Veredito (padrão auditoria)}
\begin{itemize}
  \setlength\itemsep{0.1em}
  \item \textbf{Latência:} Booleano e BM25/TF-IDF soma são ms; cosseno denso é 2~ordens mais
        lento neste corpus (122 docs $\times$ vocab) --- esperado sem índice otimizado.
  \item \textbf{Qualidade textual:} ranqueadores colocam a frase certa no topo em
        localização e fundador; AND falha (consulta longa demais); OR não ordena.
  \item \textbf{Julgamento:} pesos $0/1/2$ explicam por que $P@3$ e nDCG divergem do ``achismo''.
        Com pool parcial, $R{=}0$ (q02) impede MAP honesto --- declarar, não esconder.
  \item \textbf{Quem ``vence'':} cosseno $\approx$ BM25 no top-1 destas três; BM25 mais rápido.
        Diferença fina não é significância (Aula~16). Três consultas bastam para a demo de
        5~minutos; não bastam para coroação de modelo.
\end{itemize}

\vspace{0.4em}
\noindent\textbf{Como rodar:}
\texttt{cd estrutura/codigos} $\rightarrow$
\texttt{Rscript coletar.R} $\rightarrow$
\texttt{Rscript aula02.R} $\rightarrow$
\texttt{Rscript entrega-meio-semestre-demo.R}.

\end{document}
